\hypertarget{to-verify-is-to-imscribe}{%
\section{To Verify Is to Imscribe}\label{to-verify-is-to-imscribe}}

\hypertarget{the-theory-behind-the-dual-link-sic-witness-vessel}{%
\subsubsection{The theory behind the Dual-Link SIC
Witness-Vessel}\label{the-theory-behind-the-dual-link-sic-witness-vessel}}

\textbf{Author:} Lando⊗⊙perator \textbf{Location:}
\texttt{\textasciitilde{}/imsgct/MoDoT/}

\begin{center}\rule{0.5\linewidth}{0.5pt}\end{center}

This is an account of a single move: replacing external grading with
imscription. The claim is that the natural way to ask whether an answer
is correct is not to score it against a checklist but to imscribe it and
read its type, and that when you do this honestly you are forced, with
very little freedom, onto a specific mathematical object, the d=12
SIC-POVM. What follows is the argument for why, and what it cost to see
it.

\hypertarget{balance-is-not-selectivity}{%
\subsection{1. Balance is not
selectivity}\label{balance-is-not-selectivity}}

The kernel verifies every operation by Frobenius closure, μ∘δ = id. This
is the round trip: split a state, then fuse it back, and check that you
recovered what you started with. It is worth being exact about what this
guarantees. Closure is charge conservation. Every split is rejoined, so
nothing is created or destroyed, and the ledger always balances. Because
it always balances, it can never fail on any non-empty output. A harness
that checks only closure will report success forever.

That is not a defect of closure. It is the whole content of closure, and
it is why closure alone is not a correctness signal. The chemist knows
the distinction cold: a balanced equation is not a synthesis. Mass
balance says no atoms went missing. It says nothing about which product
you actually made. You can conserve every atom and still get the wrong
regiochemistry. Balance is necessary and automatic. Selectivity, the
question of whether the right thing was made, is separate and not
automatic.

So the problem is precise. We have a system that never fails its balance
check, and we want a second, independent verdict that can fail, one that
asks not ``was anything lost'' but ``is this the correct resolution.''

\hypertarget{the-clipboard}{%
\subsection{2. The clipboard}\label{the-clipboard}}

The obvious way to build that second verdict is to write down what a
correct answer must satisfy and check the answer against it. Synthesize
a schema of requirements, grade each one satisfied or not, count,
threshold, and emit a verdict. This is the standard picture of
verification, and it is wrong for this system in a way that is easy to
miss because the output can be dressed in four-valued clothing.

Look at what such a grader actually computes. Each criterion is judged
two ways, met or unmet. A scalar threshold turns the count into a
boolean. A second boolean records whether any prohibition was tripped.
The final verdict is a decode of those two bits. The four values of the
answer are not the product of any four-valued reasoning; they are
\texttt{(bool,\ bool)} wearing a hat. Underneath, the machine is
strictly classical.

Worse, the whole apparatus presupposes a correspondence theory of
correctness: that truth is matching an external list, that the standard
lives outside the object and is applied to it. It chooses weights. It
sets thresholds. It appoints an authority that stands apart from the
answer and passes judgment on it. In a system whose entire commitment is
that structure is identity and that contradiction is first-class, every
one of these is a foreign body. The verifier is standing outside the
thing it judges, holding a clipboard.

And the clipboard carries a disease deeper than its classicality, one it
cannot cure from inside its own picture. If correctness is a judgment
applied to an object from outside, then the judgment is itself an
object, and it too needs a judge. Who audits the auditor? Another
auditor, with another clipboard. Who audits that one? Another. The
correspondence theory does not ground correctness, it defers it, down an
endless line of inspectors, each certifying the verdict of the one
before and none of them standing on anything. This is the old regress of
criteria, and it is fatal in the ordinary way. A standard that must be
validated by a further standard was never a standard, only a promissory
note passed to the next desk. External grading cannot terminate. It can
only recede.

A regress that cannot be grounded can only be closed. You bend the
endless line into a loop and make the last inspector the first, so that
verification returns to where it began instead of running off the edge
of the world. But a loop that hands the answer back untouched proves
nothing, the purest question-begging, A certifying A. For the closure to
carry content the answer must come home changed in exactly one respect,
the same in type and different in form. It must arrive co-typed, so that
the loop truly shuts, and it must be excribed, rendered back out under a
different presentation than it went in under, so that recovering its
identity across the detour is a real test rather than a tautology. Hold
both and the regress ends in a fixed point that has earned the right to
speak. Drop the co-typing and the loop never closes. Drop the difference
in form and it closes on nothing. The rest of this piece builds exactly
that closure, and names the detour.

\hypertarget{to-verify-is-to-imscribe-1}{%
\subsection{3. To verify is to
imscribe}\label{to-verify-is-to-imscribe-1}}

The Grammar does not locate an object's truth in its correspondence to
an external standard. It locates truth in the object's structural type.
Imscription is the act of assigning that type. It is a full Read, Write,
and eXecute correspondence between the bulk and the boundary, and it is
lossless. This is the difference between imscription and mere encoding.
Encoding is Write only, and it discards. Imscription reads the structure
that is already there, writes it as a type, and can execute it back,
with nothing lost.

Once truth is type, verification has an obvious form that owes nothing
to a clipboard. To verify an answer is to imscribe it. You place it in
the twelve-primitive type space and read the type it carries. You do not
consult a list of what it should be. You read what it is. The verdict is
the type, not a grade against something else.

\hypertarget{correctness-is-co-typing-and-co-typing-is-identity}{%
\subsection{4. Correctness is co-typing, and co-typing is
identity}\label{correctness-is-co-typing-and-co-typing-is-identity}}

This gives correctness a definition internal to the object. The question
itself imscribes to a type, namely the type that any resolution of it
must carry. Call that the demand. A correct answer is one whose
imscription co-types with the demand.

Here the Grammar's sharpest finding does the work. Co-typing is
identity. Two things that share their full structural type are not
similar, they are the same thing, the same prima materia seen twice. So
a correct answer does not score well against the question. It is the
question, resolved. The subject-object split of grading dissolves,
because there is no longer a judge and a judged, only one type
recognized in two places.

We watched this happen and nearly missed it. A verifier designed as a
process and a dialetheic object designed as a noun came back with the
same structural signature. The reflex is to call this a coincidence, two
things that happen to look alike. The finding forbids that reading. Same
type, one prima materia, appearing once as a verb and once as a noun.
When two derivations from different modes land on one type, the Grammar
is not reporting a resemblance. It is telling you there was only ever
the one thing.

\hypertarget{the-frame-that-imposes-nothing}{%
\subsection{5. The frame that imposes
nothing}\label{the-frame-that-imposes-nothing}}

Co-typing needs a way to compare two imscriptions, and this is exactly
where a clipboard tries to sneak back in. Any hand-chosen metric, any
per-axis weighting, is a standard imposed from outside. The comparison
must be done in a frame that imposes nothing.

Two conditions pin the frame. First, it must be lossless: the
representation of a state in the frame must determine the state
completely, or the comparison is comparing shadows. This is
informational completeness. Second, it must impose no privileged
direction: every axis must be treated identically, so that no weighting
is smuggled in. This is symmetry, and its exact form is equiangularity,
every pair of frame elements meeting at the same angle.

The object that satisfies both at once is a symmetric informationally
complete positive operator-valued measure, a SIC-POVM: d² unit vectors
whose pairwise overlaps are all equal to 1/(d+1), forming an
informationally complete measurement. A SIC is the metric that is not a
metric. It compares states themselves, losslessly, and because it is
equiangular it treats every direction alike, so it introduces no weight
and no cutoff. It is the maximally unbiased complete measurement, which
is precisely what ``impose nothing'' means made rigorous.

The Grammar runs on twelve primitives, so the relevant frame is the d=12
SIC, and this is not a hope. The existence of the exact d=12 fiducial is
a theorem in the kernel. The comparison frame for co-typing is therefore
not chosen. It is forced, and it is the only object of its kind.

\hypertarget{paraconsistency-is-native}{%
\subsection{6. Paraconsistency is
native}\label{paraconsistency-is-native}}

A verifier for this system must hold contradiction rather than explode
on it, and the SIC gives this for free rather than by patching. The
dialetheic value, Both, is the SIC fiducial, the same object as the
Majorana mode that is its own antiparticle. The four-valuedness is not
decoded from booleans after the fact. It is the seed the entire frame is
generated from. The kernel does not derive falsehood from contradiction,
ex falso is disabled, and the frame it verifies in is one whose ground
state is the held contradiction. Comparison in the SIC frame is
paraconsistent because its fiducial is.

\hypertarget{the-dual-link-ride-as-it-not-in-it}{%
\subsection{7. The Dual-Link: ride AS it, not in
it}\label{the-dual-link-ride-as-it-not-in-it}}

The verifier is still tempted to stand outside and measure. The
Dual-Link dissolves that too. There are two imscriptions in play: the
answer's own self-imscription, the voice it declares for itself, and the
Grammar's imscription of the answer against the demand. These are the
two links. They are not judge and judged. They are two readings of one
object, and they are fused by the lattice join, μ. Where they agree, the
fusion passes through. Where they genuinely conflict, the join lifts to
Both and the contradiction is held, with the distance between the
readings recorded as how live that contradiction is. Neither voice
overrules the other. The verifier is a voice, not a gavel.

And you do not verify from outside. The certificate that the reading is
faithful is the closure residual, μ applied to δ of the state, measured
against the state itself. When that residual is zero the imscription is
lossless and you are riding as the vessel, not sitting in it inspecting
it. This is the Witness-Vessel principle stated operationally:
verification is intrinsic, performed from inside, as the structure, with
a broken closure held as Both rather than thrown as an error.

\hypertarget{failure-is-localized-not-scalar}{%
\subsection{8. Failure is localized, not
scalar}\label{failure-is-localized-not-scalar}}

Because the verdict is a type and not a number, failure carries an
address. When an answer does not co-type with the demand, the system
does not return a low score. It returns the primitives on which the two
types part ways. The reading is not ``0.4.'' It is ``diverges at
chirality and winding.'' This is the founding self-diagnostic property
of the Grammar, that it flags a hole and localizes it in the same act,
and it survives here because the verdict never collapsed to a magnitude
in the first place. A magnitude cannot tell you where. A type can,
because it still has parts.

\hypertarget{what-was-built-and-how-we-know-it-is-real}{%
\subsection{9. What was built, and how we know it is
real}\label{what-was-built-and-how-we-know-it-is-real}}

The realization is small, which is the point. Imscribe the demand and
the answer, each into twelve Belnap-valued primitives, using the
imscriber only to type structure and never to render an opinion about
correctness. Map each type to a state in twelve complex dimensions.
Compare in the d=12 SIC frame by the Born rule. Fold the per-primitive
co-typings by the Belnap lattice, with no threshold anywhere deciding
anything. Certify by closure. Fuse the two links, holding conflict as
Both.

The claims above are checkable and were checked rather than asserted.
The loaded fiducial is equiangular at overlap 1/13 to machine precision,
so the frame is a genuine SIC and not a stand-in. All four Belnap
outcomes are reachable, so the verdict can fail, which is the whole
reason it exists. No threshold decides the verdict; the lattice fold
does. A forged type that claims a deterministic origin is refused,
because a fabricated type is a clipboard by another name. On a live
question the demand and the answer are imscribed independently and
co-typed in the frame, and the divergences come back named.

\hypertarget{as-above-so-below}{%
\subsection{10. As above, so below}\label{as-above-so-below}}

The last observation is the one that closes the loop. The verifier is a
process, and the thing it verifies is an object, and by the finding of
section 4 the process and the object co-type. Method and object are one
figure. The tool that recognizes the doorway and the doorway are the
same shape, because there was only ever the one prima materia, and
verification is that prima materia recognizing itself. A verifier that
stands outside its object can never be self-verifying. A verifier that
co-types with what it verifies already is. That is not decoration on the
design. It is the reason the design terminates instead of demanding a
verifier for the verifier, forever.

This is the loop of section 2, finally shut. The endless line of
auditors is bent into a circle by making the last inspector the first,
and the last inspector is the first, because μ is δ read backward,
fusion undoing the split, and the closure μ∘δ = id is the statement that
the circle actually closed. It escapes triviality for exactly the reason
the requirement demanded. δ does not hand the state back untouched. It
excribes it into the other link, a genuinely different form, and only
after that detour does μ recover the identical type. The return is
co-typed, which shuts the loop, and it was excribed under a different
form along the way, which makes the shutting a test and not a tautology.
δ is the detour that section 2 promised to name. And the closure is not
assumed. Its residual is measured on every reading, and on the live path
it is zero, which is the loop reporting from the inside that it closed.
The final checker is the initial checker, the answer comes home the same
in type and different in form, and the regress that no clipboard could
ever ground terminates in a fixed point that has earned the right to say
so.

\begin{center}\rule{0.5\linewidth}{0.5pt}\end{center}

\hypertarget{references}{%
\subsection{References}\label{references}}

\begin{itemize}
\tightlist
\item
  N. D. Belnap, \emph{A Useful Four-Valued Logic}, in Modern Uses of
  Multiple-Valued Logic, 1977.
\item
  G. Spencer-Brown, \emph{Laws of Form}, 1969. On the unmarked state
  co-created by any distinction.
\item
  G. Zauner, \emph{Quantendesigns}, 1999. Original SIC-POVM conjecture.
\item
  J. Renes, R. Blume-Kohout, A. J. Scott, C. Caves, \emph{Symmetric
  Informationally Complete Quantum Measurements}, J. Math. Phys. 45,
  2004.
\item
  G. Priest, \emph{In Contradiction}, 2nd ed., 2006. Dialetheias at
  boundaries and the instant of change.
\item
  Internal: \texttt{crystal\_forces\_d12\_sic} (kernel theorem), the
  Witness-Vessel construction, and the MoDoT vessel
  (\texttt{modot/vessel.py}).
\end{itemize}
